% !TEX program = lualatex
% badness skip-file
\documentclass[a4paper,article,openany,12pt]{memoir}

\usepackage{preamble}

\usepackage{forest}
\forestset{
  dir tree/.style={
    for tree={
      font=\ttfamily,
      grow'=0,                 % Grow horizontally
      child anchor=west,       % Anchor children on the left
      parent anchor=south west,% Align lines properly
      anchor=west,
      calign=first,
      inner xsep=7pt,         % Spacing around text
      edge path={
        \noexpand\path [draw, \forestoption{edge}] 
        (!u.south west) +(7.5pt,0) |- (.child anchor) \forestoption{edge label};
      },
      before typesetting nodes={
        if n=1 {
          insert before={[,phantom,coordinate]}
        }{}
      },
      fit=band,
      before computing xy={l=20pt}
    }
  }
}

\usepackage{metalogo}
\NewDocumentCommand{\texlogo}{}{{\ncmfont\mdseries\TeX\ }}
\NewDocumentCommand{\lualatexlogo}{}{{\ncmfont\upshape\mdseries\LuaLaTeX\ }}
\NewDocumentCommand{\luatexlogo}{}{{\ncmfont\upshape\mdseries\LuaTeX\ }}
\NewDocumentCommand{\latexlogo}{}{{\ncmfont\mdseries\LaTeX\ }}
\NewDocumentCommand{\xelatexlogo}{}{{\ncmfont\upshape\mdseries\XeLaTeX\ }}
\NewDocumentCommand{\pdftexlogo}{}{{\ncmfont\mdseries{pdf}\TeX\ }}

\ThaiAddWords{กลุ่มตัวอักษร,กลุ่มตัวอักษรไทย,ไม่พึงประสงค์,แต้มเสีย,โหนดระยะช่องไฟ,โหนดอักขระ,โหนดจุดเลือกตัดคำ,ค่าโดยปริยาย}
\ThaiAddHyphenations{ตำ-แหน่ง,เหมือน,อีก-รอบ,ขั้น-ตอน-วิธี}

\pretolerance=1000
\tolerance=5000
\hyphenpenalty=500
\doublehyphendemerits=2000

\setlength{\parfillskip}{0pt plus 0.75\textwidth}

\begin{document}

\title{\texttt{thai-segmentation} \\ ชุดคำสั่งตัดคำภาษาไทยใน \lualatexlogo}
\author{นายภูริพัฒน์ ถัมภ์บรรฑุ}
\date{พิมพ์ฉบับแรกวันที่ : 31 ตุลาคม พ.ศ.\ 2569 \\ ฉบับปัจจุบัน : \today}
\maketitle

\begin{abstract}
  \setSpacing{1.5}
  ชุดคำสั่งนี้ (package) ใช้ตัดคำภาษาไทยใน \lualatexlogo โดยใช้ขั้นตอนวิธี \texttt{newmm} จากโปรแกรม PyThaiNLP \parencite{pythainlp-proceedings}
  ซึ่งเป็นขั้นตอนวิธีตัดคำภาษาไทยแบบใช้พจนานุกรมคำไทย และ หาตำแหน่งแยกคำด้วยยัติภังค์แบบใช้พจนานุกรมพยางค์
  ผู้ใช้งานสามารถเพิ่มคำหรือลบคำออกจากพจนานุกรมทั้งสองได้
  ระบุตำแหน่งที่จะแยกคำด้วยยัติภังค์ได้
  และสามารถควบคุมขนาดของช่องว่างระหว่างคำหรือระหว่างกลุ่มตัวอักษรในประโยคได้
\end{abstract}

\tableofcontents*

\chapter{วิธีการใช้งานและการตั้งค่าเบื้องต้น}

\section{ความต้องการขั้นต่ำของชุดคำสั่ง}

\begin{enumerate}
  \item \lualatexlogo รุ่น (เวอร์ชัน) หลังเดือนตุลาคม ปี ค.ศ. 2025 (พ.ศ. 2568)
  \item ชุดคำสั่ง \texttt{fontspec} และ \texttt{polyglossia} ของ \latexlogo
  \item ชุดแบบอักษรไทย (Thai font)
\end{enumerate}

\paragraph{หมายเหตุ :}

ณ ขณะนี้ ยังไม่สามารถใช้ชุดคำสั่งนี้กับชุดคำสั่ง \texttt{babel} ได้

\section{วิธีการดาวน์โหลดชุดคำสั่งจากแหล่งต้นทาง}

\paragraph{ในระบบปฎิบัติการ Linux :}

ให้รันคำสั่งต่อไปนี้ใน bash

\begin{bashsource}{}
mkdir -p $(kpsewhich -var-value TEXMFHOME)/tex/latex
cd $(kpsewhich -var-value TEXMFHOME)/tex/latex
git clone https://github.com/AtomThum/thai-segmentation
rm -rf .git
\end{bashsource}

สามารถตรวจสอบว่าชุดคำสั่งพร้อมใช้งานได้ด้วย
\begin{bashsource}{}
kpsewhich thai-segmentation.sty
\end{bashsource}

\paragraph{ในระบบปฎิบัติการ Windows :}

\begin{enumerate}
  \item ถ้าไม่ได้ตั้งค่า \texttt{TEXMFHOME} ไว้ ให้สร้างโฟลเดอร์ \\
        \texttt{\%USERPROFILE\%\textbackslash texmf\textbackslash
          tex\textbackslash latex\textbackslash thai-segmentation}
  \item ไปที่ยูอาร์แอล \url{https://github.com/AtomThum/thai-segmentation}
  \item กดที่ปุ่ม \texttt{Code} สีเขียว แล้วกดดาวน์โหลด ZIP ไฟล์
  \item แตกไฟล์ \texttt{thai-segmentation-main.zip}
  \item คัดลอกไฟล์ในโฟลเดอร์ที่แตกออกมาไปที่ \\
        \texttt{\%USERPROFILE\%\textbackslash tex\textbackslash
          latex\textbackslash thai-segmentation}
\end{enumerate}

\paragraph{ในระบบปฎิบัติการ Macs :}

ยังไม่มีข้อมูลเนื่องจากผู้พัฒนาชุดคำสั่งไม่มีเงินเพื่อซื้อ Macbook

\section{การเรียกใช้งานชุดคำสั่งและการตั้งค่าเอกสาร}
\label{sec:configuration}

ตั้งเอกสารให้ใช้ \lualatexlogo และตั้งให้ \latexlogo ใช้ภาษาไทยโดยใช้
\texttt{polyglossia} เลือกชุดแบบอักษรที่ต้องการโดยใช \texttt{fontspec}
ตัวอย่างเช่น

\begin{texsource}{}
% !TEX program = lualatex
\documentclass{article}

\usepackage{fontspec}
\usepackage{polyglossia}

\newfontfamily{\thaifont}{Sarabun}[Script=Thai]
\newfontfamily{\thaifontsf}{Athiti}[Script=Thai,Scale=MatchUppercase]
\newfontfamily{\englishfont}{Sarabun}[Ligatures=TeX]

\setdefaultlanguage{thai}
\setmainlanguage{thai}
\setotherlanguage{english}
\end{texsource}

ณ ขณะนี้ ยังไม่สามารถใช้ชุดคำสั่งนี้กับชุดคำสั่ง \texttt{babel} ได้

สามารถเรียกใช้ชุดคำสั่งจาก \texttt{thai-segmentation} ได้โดย

\begin{texsource}{}
  \usepackage[<options>]{thai-segmentation}
\end{texsource}
สามารถปรับค่าต่าง ๆ ใน \texttt{<options>} ได้ดังนี้
\begin{description}[font=\ttfamily]
  \item[thai-segmentation] (\texttt{true}/\texttt{false}) เปิด-ปิดการตัดคำไทย
  \item[interword-expansion] (\texttt{true}/\texttt{false})
        เปิด-ปิดการขยายช่องว่างระหว่างคำในประโยค
        (เพิ่มเติมใน\zcref{sec:interword-expansion})
  \item[interword-stretch] (ทศนิยม) จะมีผลก็ต่อเมื่อ
        \texttt{interword-expansion=true} เท่านั้น
        กำหนดขนาดที่ช่องว่างระหว่างคำสามารถขยายได้ มีหน่วยเป็นพอยต์ (pt)
        มีค่าโดยปริยายเป็น \texttt{1} ไม่ควรตั้งค่าเกินกว่านี้
        เนื่องจากอาจทำให้คำในประโยคห่างกันมากเกินไป
  \item[interword-shrink] (ทศนิยม) จะมีผลก็ต่อเมื่อ
        \texttt{interword-expansion=true} เท่านั้น
        กำหนดขนาดสูงสุดที่ช่องหว่างระหว่างคำสามารถหดได้ มีหน่วยเป็นพอยต์ (pt)
        มีค่าโดยปริยายเป็น \texttt{0.5}
        ไม่ควรตั้งค่าเกินกว่านี้เนื่องจากอาจทำให้คำชิดกันเกินไปหรือซ้อนทับกัน
  \item[intercluster-stretch] (ทศนิยม)
        กำหนดระยะห่างเพิ่มเติมระหว่างกลุ่มตัวอักษรไทย (Thai character clusters)
        จะมีผลก็ต่อเมื่อตั้งเป็นค่าทศนิยมบวก มีค่าปริยายเป็น \texttt{0}
        ควรตั้งค่าให้น้อยกว่า \texttt{interword-stretch}
  \item[dictionary] ระบุที่อยู่ (path) ของไฟล์คำที่จะใส่เพิ่มเข้าพจนานุกรมที่
        \texttt{newmm} ใช้ตัดคำ ตัวอย่างเช่น
        ถ้าโครงสร้างโฟลเดอร์เป็นแบบ\zcref{fig:folder-example} โดยไฟล์
        \texttt{file.txt} มีเนื้อหาดังนี้
        \begin{texsource}{file.txt}
ไอน์สไตน์
พลังค์
แมกซ์เวลล์
\end{texsource}
        สามารถเพิ่มคำจากไฟล์ \texttt{file.txt} เข้าพจนานุกรมได้โดยตั้งค่า \\
        \texttt{dictionaries=file.txt}
        (เพิ่มเติมใน\zcref{sec:interword-expansion})
        \begin{figure}[h]
          \centering
          \begin{forest}
            dir tree [project/ [file.txt] [main.tex] [main.pdf] ]
          \end{forest}
          \caption{ตัวอย่างโครงสร้างโฟลเดอร์}
          \label{fig:folder-example}
        \end{figure}
  \item[tcc-min-edge] (จำนวนเต็ม) หากคำใดมีตัวอักษรน้อยกว่าค่านี้
        จะไม่ใส่จุดตัดคำด้วยยัติภังค์ มีค่าโดยปริยายเป็น \texttt{2}
        (เพิ่มเติมใน\zcref{sec:interword-expansion})
  \item[show-hyph] (\texttt{true}/\texttt{false})
        เป็นเครื่องมือสำหรับผู้พัฒนาเพื่อตรวจสอบจุดตัดคำและตำแหน่งการใส่ยัติภังค์
        ไม่ควรเปิดใช้ในเอกสารฉบับสมบูรณ์
\end{description}

ตัวอย่างการตั้งค่าและเรียกใช้ชุดคำสั่ง เช่น
\begin{texsource}{}
\usepackage[
  interword-expansion=true,
  interword-stretch=0.5,
  interword-shrink=0.2,
  intercluster-stretch=0,
  dictionary=dict.txt,
  tcc-min-edge=3,
]{thai-segmentation}
\end{texsource}

\section{รายชื่อคำสั่ง}

\begin{description}[font=\ttfamily]
  \item[\textbackslash ThaiAddWords\{<คำที่หนึ่ง,คำที่สอง,...>\}] เพิ่มคำเข้าพจนานุกรมของ
        \texttt{newmm} รับค่าเหมือน \texttt{dictionaries} ในการตั้งค่าชุดคำสั่ง
        ตัวอย่างวิธีการใช้เช่น \\
        \texttt{\textbackslash ThaiAddwords\{ลากรานจ์,แฮมิลตัน,แฮมิลโทเนียน\}}
  \item[\textbackslash ThaiRemoveWords\{<คำที่หนึ่ง,คำที่สอง,...>\}] ลบคำจากพจนานุกรมของ
        \texttt{newmm} รับค่าเหมือน \texttt{dictionaries} ในการตั้งค่าชุดคำสั่ง
        ตัวอย่างการใช้เช่น \\
        \texttt{\textbackslash ThaiRemoveWords\{ยั่งยืน,อย่างง่าย\}}
  \item[\textbackslash ThaiAddHyphenation\{<คำ-ที่-หนึ่ง,คำ-ที่-สอง,...>\}] กำหนดตำแหน่งการใส่ยัติภังค์
        กรณีที่ระบบตัดคำใส่จุดตัดคำที่ไม่ตรงกับการออกเสียงหรือคำอ่านธรรมชาติ
        เช่น หากระบบตัดคำว่า "ลากรานจ์" เป็น "ลาก-รานจ์"
        ซึ่งไม่ตรงกับคำอ่านธรรมชาติ สามารถแก้ไขให้ถูกต้องได้ด้วยคำสั่ง \\
        \texttt{\textbackslash ThaiAddHyphenation\{ลา-กรานจ์\}}
  \item[\textbackslash ThaiRemoveHyphenation\{คำ-ที่-หนึ่ง,คำ-ที่-สอง\}] ใช้เหมือน
        \texttt{\textbackslash ThaiAddHyphenation}
        แต่ใช้เพื่อลบคำสั่งบังคับยัติภังค์ที่กำหนดไว้ก่อนหน้า
  \item[\textbackslash ThaiAddWordsFromFile\{<ไฟล์พจนานุกรม>\}] นำเข้าไฟล์พจนานุกรมเพิ่มเติม
        การทำงานเหมือนการกำหนดผ่านการตั้งค่า \texttt{dictionary}
        (\zcref{sec:configuration})
  \item[\textbackslash ThaiAddSyllables\{<กลุ่มอักษรที่หนึ่ง,กลุ่มอักษรที่สอง,...>\}] เพิ่มกลุ่มอักษรเพื่อใช้ในการใส่ยัติภังค์
  \item[\textbackslash ThaiRemoveSyllables\{<กลุ่มอักษรที่หนึ่ง,กลุ่มอักษรที่สอง>\}] ลบกลุ่มอักษรออกจากระบบการใส่ยัติภังค์
\end{description}

\chapter{การตัดคำและการจัดข้อความให้เต็มบรรทัด}
\label{sec:interword-expansion}

เนื่องจากภาษาไทยไม่เขียนเว้นวรรคระหว่างคำในประโยค
จึงต้องหาว่าคำคำหนึ่งเริ่มตรงไหนจบตรงไหน นั่นคือ
ต้องหาขอบเขตของคำเพื่อจะได้ตัดบรรทัดถูกและสามารถจัดข้อความให้เต็มบรรทัดได้ (text
justification) ทาง PyThaiNLP ได้พัฒนาขั้นตอนวิธี \texttt{newmm}
\parencite{pythainlp-newmm} เพื่อแก้ปัญหาการหาขอบเขตคำ โดยคร่าว \texttt{newmm}
จะตัดประโยคให้กลายเป็นกลุ่มตัวอักษรไทย (Thai character clusters)
ซึ่งอาจแทนกลุ่มอักษรควบกล้ำ คู่พยัญชนะ-สระ คู่พยัญชนะ-สระ-วรรณยุกต์ ฯลฯ
แล้วนำกลุ่มตัวอักษรต่าง ๆ ในประโยคมาเทียบกับพจนานุกรมคำไทย
แล้วค่อย\textbf{หาจุดตัดคำที่ทำให้คำทั้งหมดในประโยคยาวที่สุดเท่าที่จะเป็นไปได้}
เมื่อรู้ขอบเขตระหว่างคำแล้ว จึงสามารถตัดข้อความให้ไม่เกินความยาวบรรทัด
แล้วยืดช่องว่างหรือช่องเว้นวรรคเพื่อให้ข้อความเต็มบรรทัด แต่บ่อยครั้ง
ถึงรู้ขอบเขตคำก็อาจไม่สามารถทำให้ข้อความเต็มบรรทัดได้ เนื่องจาก
\begin{itemize}
  \item คำไทยบางคำนั้นยาวมาก เช่น \enquote{โซเดียมไฮโดรเจนกลูทาเมต} หรือ
        \enquote{สำนักงานคณะกรรมการป้องกันและปราบปรามการทุจริตและประพฤติมิชอบในวงราชการ}
        \hbox{ทำให้}หาจุดตัดคำยาก\footnote{ทั้งสองคำนี้
        อยู่ในพจนานุกรมดั้งเดิมของ \texttt{newmm}}
  \item ภาษาไทยนั้นเว้นวรรคแค่ระหว่างประโยค จึงไม่ค่อยมีช่องว่างให้ยืดด้วยซ้ำ
\end{itemize}
นอกจากหาขอบเขตคำ ชุดคำสั่งนี้จึงหาตำแหน่งแยกคำยาว ๆ
ด้วยยัติภังค์และเพิ่มช่องว่างระหว่างกลุ่มอักษรกับระหว่างคำด้วย
เพื่อการจัดเรียงตัวอักษรที่สวยงามกว่าเดิม

\section{การแยกคำด้วยยัติภังค์}
\label{sec:hyphenation}

หลังจากหาขอบเขตของคำแล้ว
ชุดคำสั่งจะไล่ดูคำอีกรอบเพื่อหาตำแหน่งแยกคำด้วยยัติภังค์
ถ้าคำไหนมีจำนวนอักษรเกิน \texttt{5} ตัวอักษร
จะนำคำไปหากลุ่มอักษรใหม่แล้วใช้ขั้นตอนวิธี \texttt{newmm}
ตัดคำให้กลายเป็นพยางค์โดยใช้พจนานุกรมพยางค์แทนพจนานุกรมคำไทย
เพื่อให้ตำแหน่งแยกคำด้วยยัติภังค์ตรงกับหลักการใช้ยัติภังค์ของราชบัณฑิตฯ
\parencite{hyphenation-orst}

แต่บางครั้ง ตำแหน่งแยกคำนั้นอาจอยู่ติดขอบคำเกิน เช่น \hbox{"\texttt{สมา-ธิ}"}
หรือ \hbox{"\texttt{ร-ยางค์}"}
จึงกำหนดเงื่อนไขไว้ว่าจะยอมให้ใส่ยัติภังค์ก็ต่อเมื่อตำแหน่งแยกคำด้วยยัติภังค์อยู่ห่างจากขอบคำทั้งสองฝั่งอยู่
\texttt{tcc-min-edge} ตัวอักษรเท่านั้น ปกติมีค่าเป็น \texttt{2}

จากที่กล่าวไว้ใน\zcref{sec:introduction} บางทีจุดตัดคำจาก \texttt{newmm}
นั้นอาจไม่ตรงกับการออกเสียงโดยธรรมชาติ สามารถแก้แบบโดยตรงโดยใช้คำสั่ง
\texttt{\textbackslash ThaiAddHyphenations}
หรือสามารถเพิ่มพยางค์เข้าพจนานุกรมพยางค์โดยใช้คำสั่ง
\texttt{\textbackslash ThaiAddSyllables} ได้เช่นกัน

ตัวอย่างคำว่า "\texttt{ลากรานจ์}" ซึ่งไม่ได้อยู่ในพจนานุกรมแต่แรก
เนื่องจากไม่ใช่คำไทยที่ใช้กันกว้างขว้าง

\begin{itemize}
  \item ถ้าใช้พจนานุกรมดั้งเดิมของ \texttt{newmm} คำว่า
        \enquote{\texttt{ลากรานจ์}} จะโดนแยกเป็นสองคำ นั่นคือ
        \enquote{\texttt{ลาก}} กับ \enquote{\texttt{รานจ์}} เนื่องจาก
        \texttt{ลาก} เป็นคำไทยที่ยาวที่สุดในพจนานุกรมคำของ
        \texttt{newmm}
  \item \ThaiAddWords{ลากรานจ์} ถ้าเพิ่มคำว่า \enquote{\texttt{ลากรานจ์}}
        เข้าพจนานุกรมคำไทย \texttt{newmm} จะเห็นคำเป็นหนึ่งคำ
        และชุดคำสั่งจะนำคำไปหาจุดตัดด้วยยัติภังค์ต่อเนื่องจากคำมีจำนวนตัวอักษรมากกว่า
        \texttt{5} แต่ \texttt{newmm} ก็ยังคงลงตำแหน่งแยกคำด้วยยัติภังค์เป็น
        \enquote{\texttt{ลาก-รานจ์}} อยู่ดี เนื่องจาก \enquote{\texttt{ลาก}}
        อยู่ในพจนานุกรมพยางค์ แต่ \enquote{\texttt{กรานจ์}} กับ
        \enquote{\texttt{รานจ์}} ไม่อยู่
  \item ถ้าเพิ่ม \enquote{\texttt{ลา}} เข้าพจนานุกรมพยางค์
        จะไม่ส่งผลเนื่องจากพยางค์ \enquote{\texttt{ลาก}} นั้นยาวกว่า
        \enquote{\texttt{ลา}} \texttt{newmm} จึงเลือก \enquote{\texttt{ลาก}}
        ซึ่งเป็นพยางค์ที่ยาวที่สุด
  \item ถ้าเพิ่ม \enquote{\texttt{กรานจ์}} เข้าพจนานุกรมพยางค์ จะทำให้
        \enquote{\texttt{newmm}} ลงตำแหน่งแยกคำด้วยยัติภังค์ถูก เนื่องจากพยางค์
        \enquote{\texttt{กรานจ์}} ยาวกว่าทั้ง \enquote{\texttt{ลา}} และ
        \enquote{\texttt{ลาก}} \texttt{newmm} จึงเลือกแยกคำเป็น
        \enquote{\texttt{ลา-กรานจ์}}
\end{itemize}

\section{การยืดช่องว่างระหว่างคำและพยางค์}
\label{sec:interword-stretch}

จะสังเกตว่าในเอกสารนี้จะมีช่องว่างเล็ก ๆ ระหว่างขอบเขตของคำ
และมีช่องว่างที่เล็กกว่าระหว่างขอบเขตพยางค์ด้วย
สามารถเปิดปิดการยืดช่องว่างระหว่างคำได้โดยตั้งค่า \texttt{interword-expansion}
กำหนดขนาดช่องว่างได้โดย \texttt{interword-stretch} และ \texttt{interword-shrink}
ดังที่กล่าวไว้ใน\zcref{sec:introduction}

\section{
  รายละเอียดเชิงลึกของวิธีการจัดบรรทัดของ \texlogo และข้อจำกัดของชุดคำสั่ง
  \texttt{thai-segmentation}
}

ในบทนี้ จะกล่าวถึงข้อจำกัดและรายละเอียดเชิงลึกของ \texlogo
ว่าด้วยระบบแต้มที่ใช้คำนวณตำแหน่งที่เหมาะสมที่สุดที่จะตัดบรรทัด
การยืดหรือหดช่องว่าง ฯลฯ แต่จะไม่ครอบคลุมเรื่องขั้นตอนวิธี Knuth-Plass

สามารถหาข้อมูลเพิ่มเติมเกี่ยวกับขั้นตอนวิธี Knuth-Plass ได้ที่บทความ
\citetitle{knuth-plass} \parencite{knuth-plass}
และสามารถหาข้อมูลเกี่ยวกับการทำงานเบื้องลึกของ \texlogo ได้ที่หนังสือ
\citetitle{texbook} \parencite{texbook} และหนังสือ \citetitle{texbytopic}
\parencite{texbytopic}

\texlogo แทนหนึ่งหน่วยข้อมูลด้วยโหนด (node) ในรายการโยง (linked list)
โหนดที่เกี่ยวข้องกับการตัดคำนั้น นั้นมีอยู่สามชนิด
\begin{description}
  \item[1. โหนดระยะช่องไฟ (glue node)] แทน
        ช่องว่างที่สามารถยืดหรือหดจากระยะปกติที่ตั้งค่าไว้ได้
        มีความสามารถในการยืดหรือหดตามที่ผู้ใช้งานกำหนด
        และทำหน้าที่เป็นจุดตัดบรรทัดแบบไม่ใส่ยัติภังค์ \texlogo
        โหนดระยะช่องไฟสามารถหดได้ถึงแค่ค่าที่ผู้ใช้งานกำหนดเท่านั้น
        \emph{แต่สามารถยืดได้ไม่จำกัด} อย่างไรก็ตาม \texlogo
        จะพยายามไม่ให้โหนดระยะช่องไฟนี้ยืดเกินค่าที่ผู้ใช้งานกำหนดผ่านขั้นตอนวิธี
        Knuth-Plass ซึ่งจะกล่าวถึงในภายหลัง \\
        สามารถกำหนดระยะต่าง ๆ ของโหนดช่องไฟที่เกิดจากเว้นวรรคได้โดยใช้ตัวแปร
        \texttt{spaceskip}\footnote{ถ้าใช้ชุดคำสั่ง \texttt{fontspec}
        สามารถตั้งค่า \texttt{WordSpace=\{A,B,C\}} ตอนเลือกแบบอักษรได้ เช่น
        \texttt{\textbackslash newfontfamily\{\textbackslash
          thaifont\}\{Niramit\}[Script=Thai,WordSpace=\{1.75,2.25,0.75\}]}}
        \begin{texsource}{}
\spaceskip=A plus B minus C
        \end{texsource}
        โดย \texttt{A} คือระยะปกติ \texttt{B} คือความสามารถในการยืด และ
        \texttt{C} คือความสามารถในการหด มีหน่วยเป็นความยาวใด ๆ ก็ได้ หรือ
        สามารถกำหนดเป็นเท่าตัวของระยะเริ่มต้นที่ \texlogo กำหนดไว้ก็ได้
        \begin{texsource}{}
\spaceskip=x\fontdimen2 plus y\fontdimen3 minus z\fontdimen4
        \end{texsource}
  \item[2. โหนดจุดเลือกตัดคำ (discretionary node)] คือจุดที่สามารถตัดคำได้
        สามารถกำหนดว่าถ้า \texlogo
        เลือกตัดตรงจุดนี้ให้พิมพ์ตัวอักษรอะไรเพิ่มเติม สำหรับภาษาไทย
        จะมีโหนดจุดเลือกตัดคำสองแบบ นั่นคือ
        แบบที่พิมพ์ยัติภังค์ก่อนขึ้นบรรทัดใหม่ ใช้ตัดระหว่างขอบเขตคำ และ
        แบบที่ไม่พิมพ์ยัติภังค์ ใช้ตัด ณ ขอบเขตคำ
        โหนดจุดเลือกตัดคำนี้ไม่สามารถยืดหดได้เหมือนโหนดระยะช่องไฟ
  \item[3. โหนดอักขระ (glyph node)] คือตัวอักขระหนึ่งตัว รวมวรรณยุกต์
\end{description}
\texlogo
\ThaiAddHyphenations{ตาม-ด้วย}
จะนำข้อความที่ผู้ใช้งานพิมพ์มาแปลงเป็นโหนดเหล่านี้ก่อนที่จะตัดข้อความให้เต็มบรรทัดโดยแปลงตัวอักษรเป็นโหนดอักขระ
แปลงเว้นวรรคเป็นโหนดระยะช่องไฟ และ แทรกโหนดจุดเลือกตัดคำ
โหนดแต่ละโหนดจะมีข้อมูลว่าโหนดไหนตามหลังด้วย เช่น ถ้าเขียนข้อความว่า
\enquote{\texttt{ก็} \texttt{ดี}} จะถูกแปลงเป็น \enquote{โหนดอักขระ ก.~ไก่}
ตามด้วย \enquote{โหนดอักขระ ไม้ไต่คู้} ตามด้วย \enquote{โหนดระยะช่องไฟ} ตามด้วย
\enquote{โหนดอักขระ ด.~เด็ก} และตามด้วย \hbox{\enquote{โหนดอักขระ สระอี}} ซึ่ง
\texlogo จะประมวลผลต่อโดยใช้ขั้นตอนวิธี Knuth-Plass
เพื่อจัดบรรทัดให้เต็มโดยอาศัยการยืดและหดของจุดยืด

กล่าวคือ ขั้นตอนวิธี Knuth-Plass จะนำโหนดต่าง ๆ ในบรรทัดมาคำนวณแต้มเสีย
(\texttt{demerit}) ซึ่งอาจจะได้แต้มมาจากการยืดโหนดระยะช่องไฟมากเกินไป
หรือการตัดคำโดยใช้ยัติภังค์ ซึ่งจะกล่าวถึงในภายหลัง ถ้าแต้มเสียเยอะ
ก็ควรเปลี่ยนไปเลือกจุดอื่นตัดบรรทัดแทน ขั้นตอนวิธี Knuth-Plass
จะพยายามทำให้แต้มเสียนี้ต่ำที่สุด
ซึ่งผู้ใช้งานสามารถปรับน้ำหนักการให้แต้มเสียของแต่ละกรณีได้ มีอยู่สองกรณี ดังนี้

\paragraph{\texttt{badness} :}

คือแต้มที่คำนวณจากการยืดหรือหดของโหนดระยะช่องไฟ
ซึ่งขึ้นกับอัตราส่วนระหว่างความสามารถในการยืดโหนดระยะช่องไฟที่ผู้ใช้ตั้งค่าไว้
(max stretch) กับระยะยืดที่ \texlogo ต้องการเพื่อให้จัดข้อความให้เต็มบรรทัดได้
(actual stretch) แต่จะไม่ให้แต้มความบิดเบี้ยวนั้นเกิน \(10000\) มีสูตรดังนี้
\begin{equation}
  \textrm{badness} = \left\lfloor 100\left( \frac{\textrm{actual stretch}}{\textrm{max stretch}} \right)^3 \right\rfloor
\end{equation}

ผู้ใช้งานไม่สามารถตั้งสูตรคำนวณแต้มความบิดเบี้ยวได้
แต่สามารถตั้งความสามารถในการยืด (max stretch) ผ่าน \texttt{spaceskip} ได้

สังเกตว่า \texttt{badness} นั้นขึ้นกับอัตราส่วนกำลังสาม ดังนั้น
ถ้าโหนดระยะช่องไฟยืดหรือหดไม่เกินความสามารถในการยืดหด แต้ม \texttt{badness}
จะน้อยมาก ๆ เนื่องจากอัตราส่วนนั้นน้อยกว่าหนึ่ง
พอนำไปกำลังสามจะได้เลขที่น้อยกว่าเดิม แต่ \texttt{badness}
จะมีผลถ้ายืดโหนดระยะช่องไฟเกินความสามารถ

\paragraph{\texttt{hyphenpenalty/exhyphenpenalty} :}

โหนดจุดเลือกตัดคำทั้งหมดจะมีแต้ม \texttt{penalty} อยู่ ซึ่งจะนำมาคำนวณเป็น
\texttt{penalty} ของย่อหน้าก็ต่อเมื่อ \texlogo
เลือกตัดบรรทัดที่โหนดจุดเลือกตัดคำเท่านั้น ถ้าแต้ม \texttt{penalty} สูง \texlogo
ก็จะพยายามไม่เลือกตัดบรรทัดที่โหนดนั้น แต่ถ้า \texttt{penalty} ต่ำ \texlogo
ก็จะยอมรับการตัดคำ ณ จุดนั้นมากขึ้น โดยแต้ม \texttt{penalty}
ของโหนดเลือกตัดคำมีค่าได้ตั้งแต่ \texttt{-10000} ถึง \texttt{10000}
โดยมีค่าปกติของโหนดจุดเลือกตัดคำที่เกิดขึ้นจากยัติภังค์เป็น \texttt{50}
สามารถปรับค่านี้ได้โดยใช้คำสั่ง \texttt{\textbackslash hyphenpenalty=...}
ถ้าตั้งค่านี้เป็น \texttt{10000} \texlogo
จะไม่ยอมตัดบรรทัดที่โหนดจุดเลือกตัดคำโดยยัติภังค์
แต่จะเลือกตัดบรรทัดที่โหนดระยะช่องไฟหรือโหนดจุดเลือกตัดคำโดยไม่ใช้ยัติภังค์เท่านั้น
เนื่องจากโหนดทั้งสองชนิดมีแต้ม \texttt{penalty} เป็นศูนย์

ผู้ใช้งานสามารถกำหนดโหนดจุดเลือกตัดคำในเอกสารเองได้โดยใช้คำสั่ง
\texttt{\textbackslash -} เช่น
\enquote{\texttt{แฮ\textbackslash -มิล\textbackslash -โท\textbackslash -เนียน}}
จุดตัดคำที่ผู้ใช้กำหนดเองนี้จะไม่ได้มีค่า \texttt{penalty} เท่ากับ
\texttt{hyphenpenalty} แต่จะใช้ค่า \texttt{exhyphenpenalty} แทน (ย่อมาจาก
explicit hyphenpenalty หรือแปลว่า \texttt{hyphenpenalty} แบบชัดแจ้ง)
มีค่าปริยายเป็น \texttt{50} ซึ่งน้อยกว่าค่าปกติของ \texttt{hyphenpenalty}
ผู้ใช้งานสามารถบังคับให้ \texlogo ตัดคำที่จุดที่ผู้ใช้งานเลือกเองได้โดยตั้งค่า
\texttt{exhyphenpenalty} เป็น \texttt{-10000}

\texlogo จะคำนวณแต้มเสียจากองค์ประกอบอื่น ๆ ด้วย
\begin{itemize}
  \item เพื่อที่จะกันไม่ให้หนึ่งย่อหน้าใช้จำนวนบรรทัดเกินกว่าที่จำเป็น \texlogo
        บวกแต้ม \texttt{linepenalty} สำหรับทุก ๆ บรรทัดที่ใช้ในย่อหน้า
        มีค่าโดยปริยายเป็น
        \texttt{10}
  \item ถ้าบรรทัดที่ติดกันมีแต้ม \texttt{badness} ห่างกันมาก เช่น
        บรรทัดบนยืดโหนดระยะช่องไฟเยอะ แต่บรรทัดล่างไม่ยืดเลย หรือหด \texlogo
        ก็จะเพิ่ม \texttt{adjdemerits} เข้าแต้มเสีย
  \item ถ้าบรรทัดที่อยู่ติดกันมียัติภังค์ทั้งสองบรรทัด ก็จะเพิ่มแต้ม
        \texttt{doublehyphendemerits} เข้าแต้มเสีย
  \item ถ้าบรรทัดที่อยู่ก่อนจบย่อหน้าเลือกจุดตัดบรรทัดที่โหนดเลือกตัดคำด้วยยัติภังค์
        จะเพิ่มแต้ม \texttt{finalhyphendemerits} เข้าแต้มเสีย
\end{itemize}
สามารถหาข้อมูลเรื่องการคำนวณแต้มเสียจากแต้มต่าง ๆ ได้จากหนังสือ
\citetitle{texbytopic} บทที่ 19 เรื่องการตัดบรรทัด \parencite{texbytopic}

ขั้นตอนวิธี Knuth-Plass จะพยายามทำให้แต้มเสียน้อยที่สุด
ซึ่งจะอธิบายพฤติกรรมประหลาด ๆ ในบางครั้งของ \texlogo
นั่นคือถ้าผู้ใช้งานตั้งค่าแต้มหนึ่งเยอะเกินไป \texlogo
อาจยอมทำอย่างอื่นเพื่อให้ได้แต้มเสียที่น้อยลง เช่น ถ้าผู้ใช้งานตั้ง
\texttt{interword-stretch} ไว้ต่ำ แต่พอดูเอกสารที่ได้อาจมีบางบรรทัดที่ยืดมากกว่า
\texttt{interword-stretch} อาจเกิดจากผู้ใช้งานตั้ง \texttt{hyphenpenalty}
มากเกินไป
ทำให้แต้มเสียของย่อหน้าจากการยืดโหนดระยะช่องไฟเกินความสามารถการยืดนั้นน้อยกว่าแต่เสียจากการเลือกโหนดจุดเลือกตัดคำ
สามารถแก้ได้โดยการตั้งแต้ม \texttt{hyphenpenalty} ให้ต่ำกว่าเดิม

สำหรับการสร้างเอกสารภาษาไทย พบว่า
\begin{texsource}{}
\hyphenpenalty=500
\doublehyphendemerits=1000
\finalhyphendemerits=2500
\end{texsource}
จะทำให้เอกสารสวยงามที่สุด

แม้ \texlogo
จะมีขั้นตอนวิธีตัดบรรทัดที่ละเอียดอ่อนและให้ผลดีที่สุดในบรรดาขั้นตอนการตัดบรรทัดทั้งหมด
แต่เนื่องจาก \texlogo นั้นออกแบบมาสำหรับภาษาอังกฤษ\footnote{\texlogo
เริ่มรองรับภาษาอื่น ๆ ในปี 1989 โดยชุดคำสั่ง \texttt{babel} ซึ่งมาหลัง \texlogo
12 ปี และเริ่มรองรับภาษาไทยในปี 2013} \texlogo
จึงไม่รู้วิธีหาตำแหน่งใส่โหนดจุดเลือกตัดคำในประโยคภาษาไทยด้วยซ้ำ ก่อนหน้านี้
คุณเทพพิทักษ์ การุญบุญญานันท์ (NECTEC)
ได้นำขั้นตอนวิธีหาตำแหน่งแยกคำด้วยยัติภังค์ของ \texlogo นั่นคือขั้นตอนวิธี
Knuth-Liang \parencite{knuth-liang} มาปรับใช้กับภาษาไทยผ่านชุดคำสั่ง
\texttt{utf8-hyph} \parencite{hyphenation-org} ซึ่งต่อมาได้นำไปใช้ในชุดคำสั่ง
\texttt{polyglossia} และ \texttt{babel} ด้วย อย่างไรก็ตาม วิธีดังกล่าวมีข้อจำกัด
เนื่องจากขั้นตอนวิธี Knuth-Liang ไม่ได้ใช้พจนานุกรมคำไทย
แต่ใช้วิธีการจับรูปแบบการแยกคำ (hyphenation pattern) ซ้อนกัน
เพื่อดูว่าตำแหน่งไหนควรใส่ยัติภังค์หรือไม่ ในภายหลัง คู่มือการใช้งานของ
\texttt{babel}
ถึงขั้นระบุว่าภาษาไทย\emph{ไม่ใช้ยัติภังค์เพื่อการแยกคำ}ด้วยซ้ำ
\parencite{babel-thai-page} ซึ่งไม่ตรงกับกฎการใช้ยัติภังค์ของราชบัณฑิต

ทางผู้พัฒนาชุดคำสั่งตระหนักดีว่าการหาตำแหน่งตัดคำของภาษาไทยนั้นเป็นเรื่องที่ท้าทาย
และเป็นโจทย์ที่สำคัญในวงการประมวลผลภาษาธรรมชาติมานานแล้ว (Natural Language
Processing - NLP) ผู้พัฒนาชุดคำสั่งจึงขอแสดงความขอบคุณอย่างสูงต่อทีมงาน
PyThaiNLP ที่ได้เปิดเครื่องมือประมวลผลภาษาไทยต่าง ๆ
แก่สาธารณะภายใต้สัญญาอนุญาตใช้ Apache 2.0 รวมถึงขั้นตอนวิธี \texttt{newmm} ด้วย
ซึ่งเปิดให้ผู้สนใจนำโปรแกรมไปใช้ แก้ไขดัดแปลงได้โดยอิสระ ซึ่งนำมาสู่ชุดคำสั่งนี้
ที่แปลขั้นตอนวิธี \texttt{newmm} จากภาษา Python เป็นภาษา Lua แล้วนำไปเชื่อมกับ
\lualatexlogo เพื่อตัดคำภาษาไทยได้ในที่สุด

ในบทถัดไป จะกล่าวถึงกระบวนการเชื่อม \lualatexlogo กับ \texttt{newmm}
และจะอธิบายชุดคำสั่งที่อยู่เบื้องหลังการทำงานของชุดคำสั่งนี้ทั้งหมด

\chapter{การเชื่อม \lualatexlogo กับ \texttt{newmm}}
\label{sec:implementation}

สามารถหาข้อมูลเพิ่มเติมเกี่ยวกับการทำงานของ \luatexlogo (ไม่ใช่
{\ncmfont\LuaLaTeX}) ได้จาก \citetitle{luatex-manual} \parencite{luatex-manual}

\section{ระบบคอลแบ็ก}

\section{ระบบแคช}


\appendix

\chapter{การใช้ปัญญาประดิษฐ์ในกระบวนการพัฒนา}

ผู้พัฒนาชุดคำสั่งได้พัฒนาส่วนต่อประสานระหว่าง \lualatexlogo และขั้นตอนวิธี
\texttt{newmm} เอง แต่ใช้ปัญญาประดิษฐ์เพื่อช่วยแปลขั้นตอนวิธี \texttt{newmm}
จากโปรแกรม \texttt{PyThaiNLP} ให้เป็นภาษา Lua โมเดลปัญญาประดิษฐ์ที่ใช้คือ Claude
Opus 5.5


\printbibliography[title=บรรณานุกรม]

\end{document}
